\begin{abstract}
点光源が物体の周囲を移動するとき、見え方の変化の多くは光がどこに当たるかによって生じる。
再照明可能な Gaussian Splatting の既存手法はこれをプリミティブごとに推定するが、本研究では光源の画像平面上で考える。
光源からガウシアンをラスタライズすると、アルファ合成重みは、各プリミティブがその光源画素の光束を受け止める割合を表す。
\ours は、この光束の\emph{沈着}を、深度モーメントと学習した光束特徴からなるライト空間アトラスに記録する。
学習時のガウシアンでも、レンダリング時の画素でも、任意の受光点がアトラスのピラミッドを参照することで、小さなネットワークで補正したフィルタ付きシャドウ判定と、沈着した特徴に対して線形な光輸送項を得られる。
また、幾何形状が形成される間は、表示用に符号化した放射輝度ではなく線形放射輝度上で損失を計算する必要があった。細い構造が複雑に入り組む二つのシーンを用いた対照実験では、この変更によりホールドアウトしたビューの PSNR が 6--7\,dB 向上した。
三つの公開ベンチマークの 18 シーンで、\ours は平均 PSNR 33.47\,dB を達成した。これは近年の四つのベースラインのうち最良の 31.59\,dB を上回り、すべてのベンチマークで上位二位以内に入る唯一の手法であり、学習時間は約 20 分である。
アブレーション実験は光輸送項の有効範囲も明らかにしている。半透明材質では特に重要である一方、ほとんどの他の材質では同規模の局所ネットワークが同等の結果を得る。
\end{abstract}
